\section{Method}\label{sec:method}

\subsection{Overview}\label{sec:overview}
Explain how the work was carried out. \Cref{fig:workflow} shows a diagram drawn directly in LaTeX with TikZ, so it needs no image file.

\begin{figure}[H]
  \centering
  \begin{tikzpicture}[
      node distance=8mm,
      box/.style={draw=primary, fill=primary!5, rounded corners,
                  minimum width=26mm, minimum height=9mm, align=center},
      arrow/.style={-Latex, thick, primary}]
    \node[box] (data) {Collect\\data};
    \node[box, right=of data] (analysis) {Analyse};
    \node[box, right=of analysis] (design) {Design\\solution};
    \node[box, right=of design] (report) {Write\\report};
    \draw[arrow] (data) -- (analysis);
    \draw[arrow] (analysis) -- (design);
    \draw[arrow] (design) -- (report);
  \end{tikzpicture}
  \caption{Workflow of the project}\label{fig:workflow}
\end{figure}

% A figure from a file in images/ (no extension needed):
% \begin{figure}[H]
%   \centering
%   \includegraphics[width=0.8\textwidth]{02-method/architecture}
%   \caption{System architecture}\label{fig:architecture}
% \end{figure}

\subsection{Formulas}\label{sec:formulas}
Numbered equations get a label like any other object. \Cref{eq:weighted-score} gives the weighted score used to compare options:
\begin{equation}
  S = \sum_{i=1}^{n} w_i \, s_i, \qquad \sum_{i=1}^{n} w_i = 1
  \label{eq:weighted-score}
\end{equation}
where $s_i$ is the score of criterion $i$ and $w_i$ its weight.

\subsection{Code}\label{sec:code}
Source code is shown with \texttt{lstlisting}, as in \cref{lst:score}.

\begin{lstlisting}[language=Python, caption={Computing the weighted score}, label={lst:score}]
def weighted_score(scores, weights):
    """Return the weighted sum of the criterion scores."""
    assert abs(sum(weights) - 1) < 1e-9
    return sum(w * s for w, s in zip(weights, scores))
\end{lstlisting}
